% Closure of the same arbitrary-SNC local projective reduction.
\begin{corollary}[Proved: projective growth on complete local punctures in all dimensions]
\label{D:complete-local-projective-growth-all-dimensions}
Let \((R,\mathfrak m,k)\) be a complete excellent local
\(k\)-algebra of characteristic \(p>0\), with algebraically
closed coefficient and residue field \(k\).
Assume every irreducible component of \(\operatorname{Spec}R\)
has dimension at least two. Put
\(U=\operatorname{Spec}R\setminus\{\mathfrak m\}\).
For every connected component \(V\) of \(U\), one has
\[
 \#\operatorname{Hom}_{\rm cont}
       \bigl(\pi_1(V),\operatorname{PGL}_r(\mathbb F_{p^f})\bigr)
          /\operatorname{PGL}_r(\mathbb F_{p^f})
                    \leq p^{C_R f r^2}                \tag{D.AP.1}
\]
for one constant \(C_R\), every \(f\geq1\), and every
\(r\geq1\). Normality and reducedness of \(R\) are not
required. No Brauer-vanishing or bounded-period condition
is imposed.

Fix \(\epsilon>0\). For each isomorphism type \(S\) in
a collection of nonabelian finite simple groups, choose
one embedding
\[
 S\hookrightarrow\operatorname{PGL}_{r_S}(\mathbb F_{q_S}),
 \qquad q_S=p^{f_S},\qquad |S|\geq q_S^{\epsilon r_S^2}.
\]
The quotient of \(\pi_1(V)\) formed by all distinct
simple quotient kernels of those types is topologically
finitely generated. Its multiplicity for type \(S\) is
at most \(|S|^{C_R/\epsilon}\). This is a simple-product
conclusion, not full finite generation of \(\pi_1(V)\).
\end{corollary}

\begin{proof}
First let \(R\) be a normal domain of dimension two.
It has a projective regular resolution, an isomorphism
off its closed point. Indeed the normal complete ring
satisfies condition (4) of
\citet[Theorem54.14.5, Tag0BGP]{StacksDSurfaceResolution},
whose proof gives normalized blowups at the singular
points. Excellence makes their normalizations finite,
so these maps are projective. Apply
\citet[Lemma54.15.6, Tag0BIC]{StacksDSurfaceResolution}
to the exceptional closed fibre. Its curve components
have finite normalization by excellence; the lemma
gives a further sequence of point blowups with SNC
support. These modifications remain over the closed
point. If the exceptional divisor is empty, blow up
the closed point once. The resulting nonempty
exceptional components are smooth proper curves
over the perfect field \(k\). Thus
\ref{D:SNC-general-projective-quadratic-count} applies.

Now let \(R\) be a normal domain of dimension \(d\geq2\).
Use the fixed slice of
\ref{D:local-surface-reduction}: choose \(d-2\) members
of a parameter system, with ideal \(I\), and put
\[
 C=(R/I)_{\mathrm{red}},\qquad
 W=\operatorname{Spec}C\setminus\{\mathfrak m_C\}.
 \]
Every component of \(C\) has dimension two.
For every connected finite étale cover of \(U\),
normalization gives a finite complete local normal
domain \(B\) over \(R\). The ideal \(IB\) has at most
\(d-2\) generators. The exact connectedness and
cohomological-dimension lemmas
\citep[Tags0ECR and0DX9]{C_StacksLocalSlice}
therefore make its punctured slice connected.
Consequently \(\pi_1(W)\twoheadrightarrow\pi_1(U)\).
This single fixed slice detects all original covers.

Normalize \(C\) and write its puncture as
\(W'=\coprod_i U_i\), where each \(U_i\) is the
puncture of a complete normal local surface
\(k\)-algebra. Their residue fields are \(k\).
Put \(a=\#\pi_0(W'\times_W W')\).
For \(G=\operatorname{PGL}_r(\mathbb F_q)\), the
surface bounds give at most \(q^{(\sum_i C_i)r^2}\)
tuples of \(G\)-torsor isomorphism classes on \(W'\).
Each such tuple has at most \(|G|^a\) descents to
\(W\). To see the exact fibre bound, fix one tuple
upstairs. An equivariant isomorphism between its
two pullbacks to a connected component of the
double product is determined by its value at one
geometric point, with at most \(|G|\) possible
values. Cocycle conditions and quotienting by
automorphisms only decrease the number.
Finite-proper effectiveness is
\citet[Expose IX, Theorem4.7 and its proof]{C_SGA1}.
Since \(|G|\leq q^{r^2}\), this gives
\(q^{(a+\sum_iC_i)r^2}\) classes on \(W\).
Surjectivity to \(\pi_1(U)\) makes restriction
injective on ambient projective conjugacy
classes. This proves (D.AP.1) for every normal
complete local domain of dimension at least two.

For the stated nonnormal ring first remove its
nilpotents, which does not change finite étale
torsors. Its finite normalization is a finite
disjoint union of complete normal local domains
\(A_j\), all of dimension at least two by the
hypothesis. Their residue fields remain \(k\).
Their punctures give a fixed finite surjection
\(U'\to U\). The normal-domain bounds just proved,
and the identical finite-proper torsor descent
argument on \(U'\times_U U'\), give
\[
                     |H^1(U,G)|\leq q^{C_Rr^2}.
 \]
There are finitely many connected components of
\(U\). A torsor on any one of them extends by
the trivial torsor on the others; thus its
isomorphism classes inject into \(H^1(U,G)\).
This proves (D.AP.1) exactly as stated.
No group-\(H^2\) or Brauer injectivity is used
in either descent step.

Distinct simple quotient kernels give distinct
ambient projective conjugacy classes under the
chosen embedding. Hence their multiplicity is
at most \(q_S^{C_Rr_S^2}\leq |S|^{C_R/\epsilon}\).
The exact generating-tuple and Goursat/compactness
argument in the proof of
\ref{D:Brauer-zero-selected-simple-product}
now applies without its former Brauer restriction.
For example \(d=\lceil C_R/\epsilon\rceil+3\)
suffices, by the same
\citet[Proposition2.6(ii), equations(2.1)--(2.2),
pp12--13]{LucchiniMarionTracey2020D}.
It gives the stated simple-product quotient.

The lower bound two on every component dimension
is retained deliberately: the proof does not
supply the corresponding representation count
for a one-dimensional punctured component.
There is no mixed-characteristic assertion,
no uniform index-linear hereditary constant,
and no unrestricted crown conclusion.
\end{proof}
